problem:  
Let \( (X_{\mathrm{sing}}, \omega_{\mathrm{sing}}, \Omega_{\mathrm{sing}}) \) be a projective Calabi–Yau threefold with isolated orbifold singularities locally modeled on \( \mathbb{C}^3 / G \), where \( G \subset \mathrm{SU}(3) \) is a finite subgroup.  
Let \( \pi : \widetilde{X} \to X_{\mathrm{sing}} \) be a **crepant resolution**, and let \( [\omega_t] \in H^2(\widetilde{X}, \mathbb{R}) \) be a family of Kähler classes converging to \( \pi^*[\omega_{\mathrm{sing}}] \) as \( t \to 0^+ \).  
Denote by \( g_t \) the corresponding Ricci‑flat Kähler metrics on \( (\widetilde{X}, [\omega_t]) \) given by Yau’s theorem.  
Assume that the exceptional divisors in \( \widetilde{X} \) collapse under \( g_t \), so that \( (\widetilde{X}, g_t) \to (X_{\mathrm{sing}}, g_0) \) in the Gromov–Hausdorff sense.

**Question:**  
Fix an orbifold point \( p \in X_{\mathrm{sing}} \) with local model \( \mathbb{C}^3 / G \).  
Let \( E_t \subset \widetilde{X} \) be the exceptional locus lying over \( p \).  
If we rescale the metric near \( E_t \) by the local collapsing scale (for instance, so that \( \mathrm{diam}_{t}^{-1}(E_t) \sim 1 \)), we obtain a sequence of pointed manifolds  
\[
(\widetilde{X}, \lambda_t^2 g_t, x_t), \quad x_t \in E_t, \ \lambda_t \to \infty.
\]

1. Does there exist a canonical limit space  
   \[
   (\widetilde{X}, \lambda_t^2 g_t, x_t) \longrightarrow (C(G), g_G, o)
   \]
   where \((C(G), g_G)\) is a (possibly asymptotically conical) Ricci‑flat Kähler manifold depending only on the local group \(G\), not on the choice of crepant resolution or degeneration path \([\omega_t]\)?  
2. If such a universal local model \((C(G), g_G)\) exists, is its Calabi–Yau deformation space canonically identifiable with the complex deformations of the singularity germ \((\mathbb{C}^3 / G, 0)\)?  

Equivalently: does each Calabi–Yau quotient singularity carry an **intrinsically defined Ricci‑flat asymptotic germ**, appearing as the local metric blow‑up limit of *any* collapsing crepant resolution?

---

Why is it a "good" problem:  
This problem refines a central open frontier in differential and algebraic geometry: understanding the **local metric structure** near collapsing exceptional sets in degenerating Ricci‑flat Kähler manifolds.  
It isolates a universal‑model conjecture—asking whether the blow‑up limit depends solely on the local group \(G\)—that no current theory can resolve.  

- **Why nontrivial:** Although Donaldson–Sun proved partial tangent cone uniqueness for limits of Kähler–Einstein metrics, the present setting involves *collapsing* Ricci‑flat metrics, where neither curvature bounds nor uniqueness results are known. The universality across crepant resolutions directly challenges our understanding of analytic moduli spaces of Calabi–Yau metrics.  
- **Bridges multiple disciplines:** The problem joins **metric limit analysis**, **Calabi–Yau birational geometry**, and **singularity theory**—asking whether analytic degenerations recognize purely algebraic equivalence classes of singularities.  
- **Precision and profundity:** It is simple to phrase—“do collapsed Ricci‑flat metrics see only the local group \(G\)?”—yet resolution would unify analytic and algebraic characterizations of Calabi–Yau singularities, possibly introducing canonical metric invariants into birational classification theory.  
- **Unresolvable with current tools:** Neither Cheeger–Colding theory nor existing AC Calabi–Yau constructions provide the universality or resolution‑independence asserted here. Proving or disproving it would require new insights into the interaction between metric convergence and complex‑algebraic moduli.  

Thus the question is mathematically sound, open, conceptually clean, and deeply integrative—a quintessential “good” research problem.